\chapter{Paul Joseph Cohen}
\label{ch1966pauljosephcohen:paul-joseph-cohen-fields-medal-1966-a-comprehensive-report}

\begin{center}\rule{0.5\linewidth}{0.5pt}\end{center}

\textbf{Prepared:} July 2026
\textbf{Subject:} Paul Joseph Cohen (b. April 2, 1934, Long Branch, New Jersey, USA; d.~March 23, 2007, Stanford, California)
\textbf{Position:} Professor, Stanford University
\textbf{Doctoral Advisor:} Antoni Zygmund (University of Chicago)
\textbf{Fields Medal Citation:} \emph{``Used technique called `forcing' to prove the independence in set theory of the axiom of choice and of the generalized continuum hypothesis.''}

\begin{center}\rule{0.5\linewidth}{0.5pt}\end{center}

\begin{prereqs}
Mathematical logic (first-order logic, model theory, completeness and compactness theorems), Zermelo--Fraenkel set theory (axioms of ZFC), cardinal and ordinal arithmetic, Boolean algebras and complete Boolean algebras.
\end{prereqs}

\begin{objectives}
Explain the method of forcing and its technical machinery including names and generic filters; prove the independence of the Continuum Hypothesis and the Axiom of Choice from ZFC \cite{cohen1963}; understand the structure of generic extensions and Cohen reals; develop familiarity with Boolean-valued models and Solovay's model.
\end{objectives}

\section{Problem History}

The Continuum Hypothesis (CH), conjectured by Georg Cantor in 1878, asserts that there is no set whose cardinality lies strictly between that of the natural numbers and the real line.

Formally, CH states $2^{\aleph_0} = \aleph_1$, where $\aleph_1$ is the smallest uncountable cardinal.

Cantor devoted the final decades of his career to attempting a proof, and the problem remained one of the most tantalizing open questions in mathematics for over half a century.

Cantor introduced the aleph notation for infinite cardinalities and developed the theory of ordinal types, but the specific gap between $\aleph_0$ and $2^{\aleph_0}$ remained impenetrable.

In 1900, David Hilbert placed CH as the first problem on his celebrated list of twenty-three problems, cementing its central status in mathematics.

Hilbert himself sketched an apparent proof, but it contained a gap that was never filled.

The difficulty of the problem suggested that CH might occupy a fundamentally different logical status from ordinary mathematical theorems---one that would require entirely new methods to resolve.

Hilbert's optimism that CH would be settled by conventional methods was echoed by many mathematicians of the early twentieth century.

However, the failure of all attempts began to suggest that CH might be independent of the standard axioms of set theory.

The Zermelo--Fraenkel axioms with the Axiom of Choice (ZFC) had been developed in the 1920s and 1930s as a foundation for mathematics, but their strength relative to CH was unknown.

The first major advance came from Kurt Gödel in 1938.

Gödel constructed the constructible universe $L$, a canonical inner model of ZFC built by transfinite recursion.

The hierarchy is defined as $L_0 = \varnothing$, $L_{\alpha+1} = \text{Def}(L_\alpha)$ (the collection of all subsets of $L_\alpha$ definable over $L_\alpha$ with parameters), and $L_\lambda = \bigcup_{\alpha < \lambda} L_\alpha$ for limit ordinals $\lambda$.

Gödel showed that $L \models \text{ZFC} + \text{CH} + \text{GCH}$, proving $\text{Con}(\text{ZFC}) \implies \text{Con}(\text{ZFC} + \text{CH})$.

The constructible universe exhibits a beautiful rigidity: every subset of every $L_\alpha$ is definable from earlier elements, leaving no room for the ``extra'' sets that would witness a failure of CH.

Gödel's technique of inner models showed that the consistency strength of CH is no greater than that of ZFC itself.

The constructible universe $L$ is the smallest inner model containing all ordinals, and it satisfies not only CH but the full Generalized Continuum Hypothesis $2^{\aleph_\alpha} = \aleph_{\alpha+1}$ for all ordinals $\alpha$.

Gödel's result established that CH cannot be disproved from ZFC (assuming ZFC is consistent).

However, the complementary question---whether CH could be proved from ZFC---remained open.

For twenty-five years, no progress was made.

The prevailing view among set theorists was that CH might eventually be settled, and many considered it probable that CH would be proved true.

Some hoped that new axioms of infinity---large cardinal axioms---might eventually resolve the question one way or the other.

The difficulty lay in the fact that existing tools for constructing models of set theory were inherently ``inner'': they started with a given model and extracted a smaller, canonical submodel.

To refute CH, one needed an ``outer'' technique---a way of adding new sets to a model while preserving the ZFC axioms.

No such method existed before Cohen's work.

The challenge was to extend a countable transitive model of ZFC to a larger model in which the power set of $\omega$ had been enlarged, while ensuring that all the ZFC axioms continued to hold in the extension.

Paul Joseph Cohen, trained as an analyst under Antoni Zygmund at the University of Chicago, approached the problem from an outsider's perspective.

His doctoral work on the uniqueness of trigonometric series had given him deep familiarity with real analysis and measure theory, but not with the technical apparatus of mathematical logic.

This fresh perspective proved decisive: Cohen was not constrained by the methodological assumptions that had limited the approaches of specialists in set theory.

He later described how his ignorance of the field's conventional wisdom freed him to try methods that logicians had dismissed.

Cohen's 1963 announcement of the independence of CH, followed by his detailed 1966 monograph, introduced forcing---a technique that has since become the single most important method in set theory.

The impact extended far beyond CH: forcing settled the independence of the Axiom of Choice, the Suslin hypothesis, Whitehead's problem, and dozens of other questions, establishing the paradigm that many statements of set theory are provably independent of the standard axioms.

The technique revealed that the universe of set theory is far richer than any single model can capture.

\section{How Was It Solved}

Cohen's method of forcing begins with a countable transitive model $M \models \text{ZFC}$ and constructs a generic extension $M[G]$ that satisfies ZFC but in which CH fails.

The construction involves three components: a partial order $\mathbb{P} \in M$, a generic filter $G \subseteq \mathbb{P}$, and a system of names for elements of $M[G]$.

A \textbf{partial order} $\mathbb{P}$ consists of conditions $p, q \in \mathbb{P}$ with a partial order $p \leq q$ meaning $p$ is stronger than $q$.

A \textbf{filter} $G \subseteq \mathbb{P}$ is non-empty, upward-closed, and closed under finite intersections.

The filter $G$ is \textbf{$M$-generic} if $G \cap D \neq \varnothing$ for every dense subset $D \in M$ of $\mathbb{P}$.

The genericity condition ensures that $G$ meets every collection of dense sets that exists within the ground model, making $G$ sufficiently generic that the resulting extension $M[G]$ satisfies all the ZFC axioms.

\subsection*{The Forcing Construction}

To negate CH, Cohen chose $\mathbb{P} = \text{Fn}(\omega_2 \times \omega, 2)$, the partial order of finite partial functions from $\omega_2 \times \omega$ to $\{0,1\}$, ordered by reverse inclusion.

Each condition $p$ is a finite set of pairs $((\alpha, n), v)$ with $\alpha < \omega_2$, $n < \omega$, and $v \in \{0,1\}$.

A stronger condition extends the function by specifying its values on additional pairs.

The generic filter determines a function $f_G: \omega_2 \times \omega \to \{0,1\}$, and for each $\alpha < \omega_2$, the function $c_\alpha(n) = f_G(\alpha, n)$ defines a new subset of $\omega$.

The partial order $\text{Fn}(\omega_2 \times \omega, 2)$ satisfies the \textbf{countable chain condition} (ccc): every antichain is countable.

This is proved by the $\Delta$-system lemma---any uncountable family of finite subsets of $\omega_2 \times \omega$ contains an uncountable $\Delta$-system (a subfamily where all pairwise intersections equal a fixed ``root''), and conditions sharing the same root are compatible.

The ccc ensures that $M$ and $M[G]$ have the same cardinals and ordinals, so no collapse of cardinals occurs.

The \textbf{forcing relation} $p \Vdash \varphi$ is defined within $M$ by recursion on formula complexity.

For atomic formulas, $p \Vdash \tau_1 \in \tau_2$ iff the set of conditions forcing this membership is dense below $p$.

The fundamental theorem states: $M[G] \models \varphi$ iff $\exists p \in G$ with $p \Vdash \varphi$.

This reduces the truth of statements in $M[G]$ to the combinatorial structure of $\mathbb{P}$ within $M$, allowing metamathematical questions to be converted into combinatorial questions about partial orders.

The \textbf{generic extension} $M[G]$ is built using \textbf{names}: every element of $M[G]$ is the interpretation $\tau_G$ of a name $\tau \in M$, where $\tau_G = \{\sigma_G : \exists p \in G, (\sigma, p) \in \tau\}$.

The collection of $M$-names is constructed in $M$ using the cumulative hierarchy.

In $M[G]$, each $c_\alpha$ for $\alpha < \omega_2$ is a distinct new subset of $\omega$, giving $2^{\aleph_0} \geq \aleph_2$ and hence $M[G] \models \neg\text{CH}$.

\subsection*{Names and the Interpretation Map}

To formalize the construction, one defines the class of $\mathbb{P}$-names $\text{Name}(\mathbb{P})$ by transfinite recursion: $\tau \in \text{Name}(\mathbb{P})$ iff every element of $\tau$ is a pair $(\sigma, p)$ with $\sigma \in \text{Name}(\mathbb{P})$ and $p \in \mathbb{P}$.

The \textbf{interpretation map} evaluates names at the generic filter: $\tau_G = \{\sigma_G : \exists p \in G, (\sigma, p) \in \tau\}$.

This map is surjective onto $M[G]$, and one proves by $\in$-induction that for every formula $\varphi(x_1, \ldots, x_n)$ and every $M$-name $\dot{x}_i$, there exists a condition $p$ forcing $\varphi(\dot{x}_1, \ldots, \dot{x}_n)$.

The \textbf{forcing theorem} guarantees that the relation $p \Vdash \varphi$ is definable in $M$ and satisfies:
\begin{enumerate}
\item $M[G] \models \varphi(\dot{x}_{1,G}, \ldots, \dot{x}_{n,G})$ iff $\exists p \in G$ with $p \Vdash \varphi(\dot{x}_1, \ldots, \dot{x}_n)$.
\item $p \Vdash \neg\varphi$ iff for all $q \leq p$, $q \not\Vdash \varphi$.
\item $p \Vdash \varphi \to \psi$ iff $\{q \leq p : q \Vdash \varphi \implies q \Vdash \psi\}$ is dense below $p$.
\end{enumerate}

\subsection*{Cardinal Preservation}

A crucial technical achievement of forcing is proving that cardinals are preserved.

For Cohen forcing $\text{Fn}(\omega_2 \times \omega, 2)$, the ccc property ensures that no cardinal is collapsed.

More generally, a partial order $\mathbb{P}$ preserves cardinals if the forcing relation respects ordinal comparisons.

The \textbf{$\kappa$-chain condition} generalizes ccc: $\mathbb{P}$ satisfies the $\kappa$-cc if every antichain has size less than $\kappa$.

The $\aleph_1$-cc (ccc) preserves $\aleph_1$ and all larger cardinals, while the $\aleph_2$-cc would preserve $\aleph_2$ and above.

Shelah's \textbf{proper forcing} theorem extends this analysis, characterizing those forcings that preserve $\omega_1$.

\subsection*{Independence of the Axiom of Choice}

For the independence of the Axiom of Choice, Cohen used a different forcing that adds a countable family of pairs of disjoint non-empty sets without a choice function.

The conditions are finite partial choice functions, and the generic filter produces a family with no uniform selector.

The resulting model satisfies all ZF axioms but not AC.

Cohen's technique was later refined by Solovay and others to produce models with remarkable properties such as the Lebesgue measure being defined on all subsets of $\mathbb{R}$.

\subsection*{Boolean-Valued Models}

An elegant alternative to the generic filter approach is the \textbf{Boolean-valued model} framework developed by Kripke and Scott.

Given a complete Boolean algebra $B$, one constructs the universe $V^{(B)}$ where every formula $\varphi$ receives a Boolean value $\|\varphi\| \in B$.

The atomic clause sets $\|x \in y\| = \bigvee_{(z,b) \in y} b \cdot \|z = x\|$, and logical connectives are interpreted by the Boolean operations: $\|\varphi \land \psi\| = \|\varphi\| \cdot \|\psi\|$ and $\|\neg\varphi\| = \|\varphi\|^*$.

Forcing then corresponds to evaluating $V^{(B)}$ at an ultrafilter $U$ on $B$: the quotient $V^{(B)}/U$ yields a classical model of ZFC.

When $B = \text{RO}(\mathbb{P})$ is the regular open algebra of the forcing poset, this recovers the standard forcing construction.

The Boolean-valued perspective provides cleaner proofs of many forcing theorems and connects forcing to the theory of Boolean-valued models.

Many logicians consider the Boolean-valued approach the more natural framework, as it avoids the need for countable transitive models and reveals the algebraic structure underlying forcing.

\subsection*{Solovay's Model}

Solovay's model, constructed using an inaccessible cardinal, is a striking application of forcing.

Solovay showed that if there exists an inaccessible cardinal $\kappa$, then there is a generic extension of $L$ in which every set of real numbers is Lebesgue measurable, has the Baire property, and satisfies the perfect set property.

This was remarkable because it was known that the negation of AC (in the form of the existence of a non-measurable set) could not be proved without some form of choice or large cardinal assumption.

The construction uses an iteration of forcings, starting from $L$ and collapsing an inaccessible cardinal to $\omega_1$.

The resulting model $\mathcal{M} = L[\mathbb{R}_\kappa]$ satisfies ZF plus the statement that all sets of reals are Lebesgue measurable.

Solovay's result established that the existence of non-measurable sets is consistent with ZF but requires some form of choice; it cannot be proved in ZF alone.

The \textbf{random real} model provides another striking application of forcing. By adding a random real---a real number that avoids all measure-zero Borel sets in the ground model---one obtains a generic extension where the added real is not constructible and where various cardinal characteristics take specific values. Random forcing uses Borel sets modulo measure zero as conditions, and the resulting generic object satisfies a strong probabilistic property: the probability that a condition belongs to any given Borel set is well-defined.

The \textbf{Mathias forcing} and \textbf{Hechler forcing} are additional examples of forcings used to manipulate cardinal characteristics. Mathias forcing adds a random-like infinite subset of $\omega$ with strong almost-disjointness properties, while Hechler forcing adds a dominating real that eventually exceeds every ground model function. These constructions illustrate the fine-grained control that modern forcing techniques provide over the structure of the continuum.

The model-theoretic perspective provided by forcing has revealed deep connections between large cardinal axioms and regularity properties of sets of reals.

This connection, known as the \textbf{Martin--Steel theorem}, shows that $\text{AD}^L(\mathbb{R})$ (the Axiom of Determinacy in $L(\mathbb{R})$) follows from infinitely many Woodin cardinals.

\section{Impact of This Discovery}

Cohen's independence results, combined with Gödel's consistency proof, settled Hilbert's first problem definitively: CH is independent of ZFC.

This was a watershed moment in the history of mathematics, demonstrating that the standard axioms of set theory do not determine the truth value of a natural and important mathematical statement.

The philosophical implications were profound, raising questions about mathematical truth, the role of axioms, and the limits of the axiomatic method.

Many mathematicians were disturbed by the result, while others saw it as revealing the genuine freedom inherent in mathematical foundations.

The introduction of forcing transformed set theory from a subject concerned primarily with consistency proofs into a vibrant field where independence results are the norm rather than the exception.

Within a decade, forcing had been used to prove the independence of the Suslin hypothesis, the existence of Suslin trees, the Whitehead problem in group theory, and numerous statements about large cardinals and descriptive set theory.

Each independence result opened new lines of inquiry about which additional axioms might be adopted.

The method stimulated the search for new axioms that might settle the statements left independent by ZFC.

Martin's axiom, the Proper Forcing Axiom, and the Maximality Principle represent attempts to find ``natural'' additional axioms.

Large cardinal axioms---measurable, strongly compact, and Woodin cardinals---provide a hierarchy of consistency strengths that partially orders many independence results.

The search for ``the right axioms'' for set theory remains one of the most active areas of foundational research.

Iterated forcing (Martin, Solovay) and proper forcing (Shelah) extended the method to construct models with prescribed combinations of cardinal invariants.

These techniques underpin modern combinatorial set theory, including the study of $\mathcal{P}(\omega)/\text{fin}$, the Cichoń diagram, and pcf theory.

The ability to manipulate cardinal invariants independently has revealed a landscape of consistency results of extraordinary richness.

Cohen's earlier work in harmonic analysis---the Cohen--Hewitt factorization theorem and the characterization of idempotent measures on locally compact groups---demonstrated his broad mathematical abilities.

The factorization theorem states that any element of a Banach algebra with a bounded approximate identity can be written as a product of two elements, a result of independent interest in functional analysis.

The model-theoretic perspective introduced by forcing reshaped the entire foundations of mathematics.

By treating set-theoretic universes as mathematical objects that can be compared, extended, and manipulated, forcing brought a new level of precision to foundational questions.

The method also influenced other areas of mathematics: topologists used forcing to construct counterexamples in general topology, algebraists applied it to questions about modules and vector spaces, and analysts utilized forcing techniques in the study of Banach spaces.

The philosophical debate about the nature of mathematical truth was reignited by Cohen's results.

Platonists argued that CH has a definite truth value, even if ZFC cannot determine it, while formalists saw the independence result as demonstrating the conventional nature of axiomatic systems.

The multiverse interpretation, discussed further below, offers a third perspective that takes the plurality of models seriously.

\section{Modern Views and Improvements}

The Boolean-valued approach to forcing, developed by Kripke and Scott, presents forcing through complete Boolean algebras $\text{RO}(\mathbb{P})$, where the Boolean value $\|\varphi\|$ determines whether $\varphi$ holds in $M[G]$.

This algebraic presentation provides cleaner proofs and connects forcing to the theory of Boolean-valued models.

The Boolean-valued universe $V^{(B)}$ is an alternative to the countable transitive model approach, and many logicians consider it the more natural framework.

Iterated forcing with countable support allows the construction of models by iterating the forcing through ordinal lengths.

Martin and Solovay's method, and Shelah's proper forcing iteration, produce models with specific combinations of cardinal invariants, settling many questions about the structure of $\mathcal{P}(\omega)/\text{fin}$ and the Cichoń diagram of cardinal characteristics.

The theory of iterated forcing has become a vast subject in its own right.

Proper forcing (Shelah) generalizes ccc forcing while preserving $\omega_1$.

A forcing $\mathbb{P}$ is proper if for every sufficiently large $\theta$, every countable $M \prec H(\theta)$ with $\mathbb{P} \in M$, and every $p \in \mathbb{P} \cap M$, there is a condition $q \leq p$ that is $(M, \mathbb{P})$-generic.

Proper forcing has resolved deep questions beyond the reach of ccc methods, including Martin's Maximum and the preservation of $\omega_1$ under iterations.

Set-theoretic geology (Hamkins) studies ground models and their extensions.

The \textbf{mantle}---the intersection of all grounds---is a canonical inner model.

The modal logic of forcing investigates possible worlds connected by forcing relations, raising philosophical questions about the set-theoretic multiverse.

The multiverse view, championed by Hamkins, holds that there is no single true universe of sets but rather a plurality of equally valid universes connected by forcing.

Forcing axioms continue to drive research.

The Strong Forcing Axiom and Woodin's $\mathbb{P}_{\max}$ axiom determine specific models where many independent statements are settled.

Woodin's Ultimate-$L$ program aims to construct a canonical inner model that decides CH and other central statements, potentially providing a definitive resolution to the foundational questions raised by Cohen's work.

Recent developments in inner model theory have produced increasingly sophisticated canonical models for large cardinals, with Woodin cardinals now analyzed up to rank-into-rank hypotheses.

The \textbf{Core Model} $K$, developed by Mitchell and Schimmerling, provides a canonical inner model below a measurable cardinal.

Extending $K$ to accommodate stronger large cardinals has been one of the central programs in set theory for decades.

The \textbf{Hod Conjecture} proposes that the hereditarily ordinal-definable sets $HOD$ form a canonical inner model that can accommodate large cardinals.

In the realm of forcing and determinacy, the \textbf{Martin--Solovay theorem} established that a measurable cardinal is sufficient to construct a model where all projective sets are Lebesgue measurable.

This was later strengthened by Martin and Steel to show that $n$ Woodin cardinals yield determinacy for $\Pi^1_n$ sets of reals.

The study of \textbf{cardinal invariants} of the continuum, organized in the Cichoń diagram, has been greatly refined using modern forcing techniques.

These eight cardinal characteristics---including the adding number $\mathfrak{a}$, the splitting number $\mathfrak{s}$, and the reaping number $\mathfrak{r}$---describe the minimal size of families of sets with various combinatorial properties.

The Cichoń diagram has been completely mapped, with each consistent ordering of these cardinals realized by an appropriate forcing construction.

The \textbf{proper forcing axiom} (PFA) and its variants, such as Martin's Maximum (MM), have emerged as natural candidates for additional axioms that settle many independent statements.

PFA implies the negation of CH and determines many cardinal invariants, while MM, the strongest known forcing axiom, settles even more.

The consistency strength of PFA is one inaccessible cardinal, while MM requires a Mahlo cardinal.

These forcing axioms have found applications far beyond set theory, including in topology, functional analysis, and algebra.

The model-theoretic study of forcing has also deepened our understanding of the set-theoretic universe.

Hamkins' \textbf{button} and \textbf{rigging} concepts classify statements according to their modal status across forcing extensions.

A statement is a \textbf{necessity} if it holds in all forcing extensions, a \textbf{contingency} if it holds in some but not all, and a \textbf{possibility} if it fails in some but not all.

The modal logic of forcing is precisely S4.2, a result connecting the syntactic structure of forcing to well-known modal logics.

The \textbf{optimism conjecture}, proposed by Friedman and Holy, predicts that every true statement of ZFC is forced by some poset in $L$, suggesting a deep connection between truth and provability in the constructible universe.

The \textbf{diamond principle} $\diamondsuit$, introduced by Jensen, has become a powerful tool in combinatorial set theory. $\diamondsuit$ asserts the existence of a sequence $\langle A_\alpha : \alpha < \omega_1 \rangle$ such that $A_\alpha \subseteq \alpha$ for each $\alpha$, and for every $A \subseteq \omega_1$, the set $\{\alpha < \omega_1 : A \cap \alpha = A_\alpha\}$ is stationary. This principle holds in $L$ and is independent of ZFC, providing a fine-structural tool for constructing pathological objects.

The \textbf{square principles} $\square_\kappa$, also studied by Jensen, form a hierarchy of combinatorial principles that have deep connections to the structure of inner models and to the failure of large cardinal axioms. The failure of $\square_\omega$ is implied by the existence of a measurable cardinal, linking combinatorial set theory directly to large cardinal theory.

The modern study of \textbf{inner models} for large cardinals has produced the \textbf{Mitchell--Steel core model}, which accommodates Woodin cardinals and beyond. The construction uses iterated ultrapowers and extenders to build canonical models that mirror the structure of the universe $V$ as closely as possible below the large cardinal in question.

These modern developments continue to build on Cohen's foundational insight that the set-theoretic universe is far richer than any single model can capture.

\begin{center}\rule{0.5\linewidth}{0.5pt}\end{center}

\section*{Summary}
\addcontentsline{toc}{section}{Summary}

Cohen invented the method of forcing, proving that the Continuum Hypothesis and the Axiom of Choice are independent of ZFC set theory.

This technique became the central tool of set theory and one of the most important methods in mathematical logic, resolving Hilbert's first problem and revealing the genuine independence of many mathematical statements from the standard axioms.

The forcing construction, generic filters, and Boolean-valued models have provided powerful new ways to understand the structure of set-theoretic universes.

Subsequent developments in proper forcing, forcing axioms, inner model theory, and the multiverse interpretation have expanded Cohen's legacy into a rich and active research program.

Solovay's model of all Lebesgue measurable sets and the modern study of cardinal invariants demonstrate the far-reaching impact of forcing techniques.

The philosophical implications of Cohen's work continue to provoke debate about the nature of mathematical truth and the role of axioms in mathematics.
